\chapter{Introdução}\label{cap:introducao}

A seleção do algoritmo mais adequado para um determinado problema de aprendizado de máquina (\glsadd{ml}ML, \textit{Machine Learning}) é, reconhecidamente, um dos maiores desafios enfrentados por profissionais da área. Isso ocorre porque não existe um algoritmo universalmente superior para todos os tipos de dados, conforme reforçado pelo Teorema \textit{No Free Lunch} \citeonline{Wolpert1997}, que demonstra formalmente que, em média sobre todos os problemas possíveis, qualquer dois algoritmos de otimização ou aprendizado têm desempenho equivalente — implicando que a superioridade de um algoritmo em um domínio necessariamente vem acompanhada de inferioridade em outro. As características intrínsecas dos dados — como dimensionalidade, balanceamento de classes, distribuição e presença de ruídos — impactam diretamente o desempenho dos modelos, tornando inviável a adoção de uma abordagem única e generalista.

Nesse cenário, o desenvolvimento de pipelines eficientes em ML exige uma combinação cuidadosa de três componentes fundamentais: representação, avaliação e otimização, conforme proposto por \citeonline{Domingos2012}. A representação define o espaço de hipóteses possível; a avaliação, os critérios que distinguem bons modelos de ruins; e a otimização, os métodos de busca que encontram o modelo mais adequado. Tradicionalmente, a construção desses pipelines depende fortemente da expertise humana, que realiza sucessivas tentativas e ajustes manuais.

Com o objetivo de reduzir esse esforço, surgem as abordagens de \glsadd{automl}\textit{AutoML} (\textit{Automated Machine Learning}), que automatizam desde a seleção de algoritmos até o ajuste de hiperparâmetros. Deste modo, como destacam \citeonline{Schoenfeld2020} e \citeonline{Vanschoren2018}, a eficácia dos sistemas de AutoML se amplia quando são utilizados modelos de meta-aprendizado, capazes de aprender a partir da experiência acumulada em múltiplos problemas anteriores.

O meta-aprendizado, ou aprendizado para aprender, estrutura-se na coleta de metadados e na extração de \textit{meta-características} dos conjuntos de dados — atributos que capturam aspectos estatísticos, estruturais e informacionais dos dados \citeonline{Schoenfeld2022}. Com isso, torna-se possível treinar metamodelos capazes de prever, de maneira informada, quais algoritmos ou pipelines possuem maior chance de sucesso em um novo problema.

A avaliação de pipelines precisa considerar conjuntamente o classificador e o \textit{pré-processamento dos dados}. Escalas divergentes, valores ausentes e escolhas de codificação ou seleção de atributos podem alterar o desempenho de uma configuração. Além de comparar essas combinações, é necessário estimar suas transformações somente nos dados de treino, preservando o teste para avaliar a capacidade de generalização.

Diante desse cenário, este trabalho investiga o ranqueamento e a recomendação de pipelines de classificação por meta-aprendizado. A proposta compreende: (i) auditar um acervo histórico de 201 tarefas OpenML e reconstruir os escores de 60 pipelines sobre 193 tarefas elegíveis, com pré-processamento ajustado somente no treino e amostras equivalentes para todos os classificadores; (ii) representar cada tarefa rotulada por um vetor estatístico determinístico de 55 descritores e nove zeros; e (iii) comparar k-NN e Random Forest multi-saída com baselines, usando validação por grupos, ajuste interno e comparação controlada com os escores históricos. A avaliação inclui 165 grupos e permanece retrospectiva, sem inferir superioridade universal a partir das métricas do benchmark.

A avaliação de recomendações exige distinguir representação, preditor e benchmark. A revisão de \citeonline{Baratchi2024} organiza AutoML em espaço de busca, estratégia e avaliação; essa distinção motiva controlar o arsenal ao comparar metamodelos.

\section{Objetivos}\label{sec:objetivos}

Este trabalho tem como objetivo geral desenvolver e avaliar um sistema de meta-aprendizado baseado em embeddings estatísticos de conjuntos de dados para ranqueamento e recomendação de pipelines de classificação, capaz de identificar, a partir das características do conjunto de dados de entrada, quais configurações de pré-processamento e classificação apresentam maior probabilidade de sucesso preditivo. Para alcançar esse objetivo geral, foram definidos os seguintes objetivos específicos:

\begin{itemize}
    \item Avaliar, por meio da construção de um meta-conjunto — repositório de conhecimento contendo descrições de pipelines, embeddings dos conjuntos de dados e rótulos de desempenho obtidos por validação cruzada estratificada —, se metadados de desempenho e meta-características de conjuntos de dados carregam sinal suficiente para diferenciar boas e más combinações de pré-processamento e classificação;
    \item Construir e avaliar um modelo preditivo capaz de estimar, a partir de metacaracterísticas dos dados, o ranking de adequação de diferentes configurações de pré-processamento e classificação.
\end{itemize}

\section{Justificativa}\label{sec:justificativa}

A relevância deste trabalho está na articulação entre duas decisões experimentais. A primeira é avaliar conjuntamente escalonamento, seleção de atributos e classificação, preservando a identidade de cada pipeline. Essa abordagem permite investigar como diferentes combinações se comportam no mesmo conjunto de dados e evita atribuir ao classificador, isoladamente, os resultados de todo o fluxo de processamento.

A segunda decisão refere-se à forma de avaliação: recomendar apenas um pipeline não descreve a qualidade da ordenação dos demais candidatos. O ranqueamento permite examinar tanto a associação entre as ordens real e predita quanto a utilidade das primeiras recomendações. Por isso, este estudo combina correlações de postos, NDCG e regret, sem pressupor que uma única métrica represente todos esses aspectos.

O trabalho investiga essas decisões ao: (i) avaliar pipelines como combinações de pré-processamento e classificação; (ii) adotar ranqueamento como tarefa de recomendação; e (iii) utilizar métricas complementares para avaliar ordenação e utilidade das recomendações. O acervo histórico de 201 tarefas e 60 pipelines foi auditado, e a avaliação principal usa 193 tarefas em 165 grupos com escores integralmente reconstruídos. A disponibilização de configurações, índices e resultados por fold integra a contribuição. Não se infere superioridade por escala ou métricas de protocolos distintos.


\section{Estrutura do Trabalho}\label{sec:estruturaTrabalho}

Este trabalho está organizado em cinco capítulos:

\begin{itemize}
    \item \textbf{Capítulo \ref{cap:introducao} – Introdução:} apresenta o contexto, o problema de pesquisa, os objetivos, a justificativa e a estrutura geral do trabalho;

    \item \textbf{Capítulo \ref{cap:referencialTeorico} – Referencial Teórico:} discute os fundamentos de aprendizado de máquina, AutoML, meta-aprendizado, meta-características para caracterização de conjuntos de dados, aprendizado ativo, estrutura de pipelines, componentes de pré-processamento e classificação, embeddings de conjuntos de dados e métricas de ranqueamento, e revisa os principais trabalhos relacionados nas áreas de representação de conjuntos de dados por embeddings, sistemas de recomendação de algoritmos e pipelines via meta-aprendizado, e análise de complexidade de fronteiras de decisão;

    \item \textbf{Capítulo \ref{cap:materialemetodos} – Materiais e Métodos:} descreve os materiais utilizados, a metodologia experimental em duas fases — construção do meta-conjunto e meta-aprendizado por ranqueamento —, os preditores k-NN e Random Forest, a auditoria, a validação por grupos e as métricas de ranking;

    \item \textbf{Capítulo \ref{cap:experimentos} – Experimentos e Resultados:} apresenta a reconstrução integral da Fase 1, a comparação controlada com os escores históricos, a avaliação de metamodelos por grupos e o diagnóstico de erros;

    \item \textbf{Capítulo \ref{cap:conclusoeseperspectivas} – Conclusão:} sintetiza as contribuições do trabalho, verifica o atendimento dos objetivos, discute limitações identificadas e propõe direções para pesquisas futuras.
\end{itemize}
